\chapter{Composition and measurement}

\section{The chapter where everything connects}

The previous four chapters developed each slot of the lens
independently. This one wires them together. We will do two things in
parallel.

First, the \emph{worked composite example}: a Flow-Lenia substrate
coupled to a population of LLM agents that name the creatures the
field produces. Every slot of the lens will be filled in. By the end of
the section you will have a full, simulator-ready specification of a
hybrid system, written down in the lens's vocabulary.

Second, the \emph{measurement} side. We have promised that emergence
means Hoel-style causal emergence and novelty means observer-relative
divergence. This chapter writes both down quantitatively. Effective
information, causal emergence at a coarse-graining, the $\Omega$
metric for residence-time-based open-endedness, and the relationship
between observer-side scores and substrate-side dynamics.

A reader who wanted only the worked example can read sections~8.2--8.6
and stop. A reader who wanted only the measurement framework can skip
to section~8.7.

\section{The Flow-Lenia + LLM composite}

The composite has two component substrates. Flow-Lenia is a
mass-conserving continuous CA on a 2D toroidal lattice that produces
discrete, drifting, sometimes-self-replicating creatures. An LLM
substrate carries a population of language-model agents whose state is
a prompt-conditioned policy. The interesting science of the composite
sits at the interface: the LLM agents watch the Flow-Lenia field, name
the creatures they see, and the names feed back into the field's
viability landscape.

The system is intentionally elaborate. Most of the slots can be
ablated independently --- replace the LLM agents with random tokens,
swap the autopoietic filter for a Markov-blanket filter, collapse the
multiplex topology --- and the system remains an instance of the same
lens. Slot-wise comparison across ablations is what the framework is
designed to enable.

\section{Substrate}

The substrate is the coupled composition
\[
  \mathbf{X} \;=\; \mathbf{X}_{\mathrm{FL}} \boxtimes_{\kappa}
  \mathbf{X}_{\mathrm{LLM}}.
\]
The components:

\begin{itemize}[leftmargin=1.5em,itemsep=2pt]
  \item $\mathbf{X}_{\mathrm{FL}} = ([0,1]^L,\, \Phi_{\mathrm{FL}},\,
    \mu_0^{\mathrm{FL}})$ with $L = (\mathbb{Z}/H\mathbb{Z}) \times
    (\mathbb{Z}/W\mathbb{Z})$ a toroidal lattice, $\Phi_{\mathrm{FL}}$
    the Flow-Lenia kernel-convolution-and-growth update with local
    mass conservation, and $\mu_0^{\mathrm{FL}}$ an initial seed (a
    small bump of mass in the centre of the lattice).
  \item $\mathbf{X}_{\mathrm{LLM}} = (\Sigma^{*},\, \mathrm{id},\,
    \mu_0^{\mathrm{LLM}})$ with $\Sigma^*$ the space of token
    sequences. The LLM substrate's dynamics is the identity (the LLM
    is queried, not autonomously evolved), and the initial measure
    $\mu_0^{\mathrm{LLM}}$ is some prompt prior --- a distribution
    over initial agent prompts.
\end{itemize}

The coupling datum $\kappa = (\kappa_{\mathrm{FL} \to \mathrm{LLM}},\,
\kappa_{\mathrm{LLM} \to \mathrm{FL}})$:

\begin{itemize}[leftmargin=1.5em,itemsep=2pt]
  \item $\kappa_{\mathrm{FL} \to \mathrm{LLM}}$ renders the Flow-Lenia
    frame to a text prompt the LLM consumes. ``Here is the current
    field. Name the creatures you see.''
  \item $\kappa_{\mathrm{LLM} \to \mathrm{FL}}$ supplies a
    per-creature scalar viability bonus indexed by which creatures the
    LLM has named. Named creatures get a small persistence advantage.
\end{itemize}

The coupling is asymmetric: LLM state shapes Flow-Lenia viability;
Flow-Lenia state shapes LLM prompts. Neither side directly modifies the
other's dynamics. The composition is well-typed; each component
remains identifiable.

\section{Entities}

For each connected component $S \subseteq L$ of $\{p \in L :
x^{\mathrm{FL}}_p \geq \theta\}$ above mass-threshold $\theta$, define
the scale projection
\[
  \pi_{\mathrm{FL}}\!\big((x_p)_{p \in S}\big) \;=\;
  \Big(\frac{1}{m}\sum_{p \in S} p\, x_p,\, m\Big), \quad
  m := \sum_{p \in S} x_p,
\]
giving an $\mathbb{R}^2 \times \mathbb{R}_{\geq 0}$-valued
centroid-and-mass description of each Flow-Lenia creature.

LLM agents are atomic entities: $S$ is a singleton, $\pi$ is the
identity.

\section{Variation}

The variation channel composition
\[
  V \;=\; V_{\mathrm{FL}} \diamond_\rho V_{\mathrm{LLM}}
\]
with:

\begin{itemize}[leftmargin=1.5em,itemsep=2pt]
  \item $V_{\mathrm{FL}}(x;\theta) = \mathcal{N}(x,\, \sigma^2 I)$
    applied to the parameter-localised Flow-Lenia kernel weights. Each
    creature carries its own kernel parameters; mutation is
    per-creature.
  \item $V_{\mathrm{LLM}}(x;\theta)$ resamples each token of the
    agent's prompt with probability $p_{\mathrm{rewrite}}$ from the
    LLM marginal $q(\cdot \mid \mathrm{context})$.
  \item $\rho(d\theta'_{\mathrm{LLM}} \mid x_{\mathrm{FL}},
    x_{\mathrm{LLM}}, \theta_{\mathrm{LLM}})$ re-weights the LLM token
    marginals proportionally to the recent naming-success score of
    the named creatures, biasing prompt mutations toward referents
    that paid off in the last $W$ ticks.
\end{itemize}

\section{Viability}

The viability filter $F = F^*$ is hierarchical (Chapter~6):

\begin{itemize}[leftmargin=1.5em,itemsep=2pt]
  \item Each Flow-Lenia entity $e_{\mathrm{FL}} = (S, \pi_{\mathrm{FL}})$
    passes iff its support $S$ is closed under the local
    mass-conservation reaction network --- autopoietic closure on the
    Flow-Lenia ``reaction'' of mass transport.
  \item Each LLM agent passes iff it has successfully named at least
    one Flow-Lenia entity in the last $W$ ticks --- a Minimal
    Criterion Coevolution predicate.
  \item $F^*$ accepts the joint configuration iff both component
    filters pass and $|\mathcal{E}_{t+1}^{\mathrm{FL}}| \geq 1$ (a
    non-extinction top-level predicate).
\end{itemize}

The hierarchical filter is the place where the autopoietic and MCC
formalisms meet without being merged. Each formalism applies at its
own level; the composite is accepted iff each level passes.

\section{Interaction topology}

The interaction topology is a three-layer multiplex with weights
$(\alpha_1, \alpha_2, \alpha_3)$ summing to one:

\begin{itemize}[leftmargin=1.5em,itemsep=2pt]
  \item $T_t^{(1)}$: spatial-overlap on $L$ for Flow-Lenia entities.
    Hyperedges = pairs of creatures whose supports intersect within
    radius $r$.
  \item $T_t^{(2)}$: uniformly random pairing of LLM judges and
    namers each tick.
  \item $T_t^{(3)}$: bipartite stigmergic layer connecting an LLM
    agent to each Flow-Lenia entity it has named.
\end{itemize}

The layer weights $(\alpha_1, \alpha_2, \alpha_3)$ are tuning knobs.
Increasing $\alpha_3$ amplifies the LLM-to-Lenia feedback; decreasing
it isolates the two substrates.

\section{Observer}

The observer is the pair $O = (O_\Omega,\, O_{\mathrm{CLIP}})$ with
window $w = W$:

\begin{itemize}[leftmargin=1.5em,itemsep=2pt]
  \item $O_\Omega$ is the residence-time $\Omega$ metric on the time
    series of distinct Flow-Lenia attractors.
  \item $O_{\mathrm{CLIP}}$ is the mean cosine novelty of CLIP
    embeddings of (frame, name) pairs against an archive of size $A$.
\end{itemize}

The observer carries state: the attractor history for $O_\Omega$ and
the CLIP archive for $O_{\mathrm{CLIP}}$.

\section{Iteration rule}

The system iterates on $Z = X \times \mathcal{M}_+(X)$ with
\[
  (x_{t+1},\, \mu_{t+1}) \;=\;
  \big(\widetilde S_{F^*} \circ \widetilde V_{T_t} \circ
       \widetilde\Phi\big)(x_t,\, \mu_t).
\]
The observer runs alongside, consuming a window $W$ of trajectory at
each tick:
\[
  (o_t,\, a_{t+1}) \;=\;
  O_s\!\big((x,\mu)_{t-W:t+1},\, a_t\big).
\]
That is the system. Five typed slots, one coupled-composition
datum, one variation-channel coupling rule, one hierarchical viability
filter, one multiplex topology, one stateful observer. Specified in
under a page.

\section{Measurement, part one: effective information}

We move from specification to measurement. The book's commitment is
that emergence is Hoel-style causal emergence. Here is the formula.

\subsection*{The setup}

Let $M = (X, P)$ be a discrete-time, finite-state Markov chain on a
state space $X$ of size $|X| = n$, with transition matrix $P$ of size
$n \times n$. Row $i$ of $P$ is the conditional distribution over
next states given current state $i$:
\[
  P_{ij} = \Pr[x_{t+1} = j \mid x_t = i].
\]

\subsection*{The intervention distribution}

The effective information is the mutual information between current
and next state \emph{under a particular distribution over the current
state}. The Hoel convention is to use the uniform distribution
$\mathbf{H}_{\max}$ over $X$ --- each cause is intervened upon
equiprobably, in the spirit of Pearl's notion of intervention.

This is a convention, not a derivation. Other intervention
distributions are defensible; results are sensitive to the choice.

\subsection*{The formula}

Under uniform intervention, the effective information is
\[
  \mathrm{EI}(M) \;=\; \frac{1}{n} \sum_{i=1}^{n} D_{\mathrm{KL}}\!\big(
  P_{i,\cdot} \,\|\, \bar P_\cdot\big),
\]
where $\bar P_\cdot = \frac{1}{n} \sum_i P_{i,\cdot}$ is the average
row of $P$. Equivalently,
\[
  \mathrm{EI}(M) \;=\; \mathbf{H}(\bar P_\cdot) - \frac{1}{n} \sum_i
  \mathbf{H}(P_{i,\cdot}),
\]
the difference between the entropy of the average effect distribution
and the average row entropy. The right-hand side decomposes neatly:
the first term measures how spread out the effect distribution is
(non-degeneracy), the second measures how certain each row's
prediction is (determinism). The full formula is in the math primer.

\subsection*{Causal emergence}

A coarse-graining $\Pi : X \twoheadrightarrow X_\Pi$ is a measurable
partition of $X$ into macro-states. Each micro-state $i$ belongs to a
unique macro-class $\Pi(i) \in X_\Pi$. The induced macro chain
$M^\Pi = (X_\Pi, P^\Pi)$ has transitions
\[
  P^\Pi_{IJ} \;=\; \frac{1}{|\Pi^{-1}(I)|}
  \sum_{i \in \Pi^{-1}(I)} \sum_{j \in \Pi^{-1}(J)} P_{ij}.
\]
The averaging convention here is the standard one; alternatives exist
and the literature reports results under several.

The causal emergence at coarse-graining $\Pi$ is
\[
  \mathrm{CE}(M, \Pi) \;=\; \mathrm{EI}(M^\Pi) - \mathrm{EI}(M),
\]
and $M$ is causally emergent iff $\sup_\Pi \mathrm{CE}(M, \Pi) > 0$.

\subsection*{From the framework's $\mathcal{S}$ to a Markov chain}

A substrate-agnostic system $\mathcal{S} = (\mathbf{X}, V, F, T, O)$
induces a Markov chain on the joint state space $Z = X \times
\mathcal{M}_+(X)$:
\[
  M(\mathcal{S}) \;=\; \big(Z,\, \widetilde S_F \circ \widetilde V_T
  \circ \widetilde\Phi\big).
\]
The transition matrix is the joint-update kernel. For substrate-as-
population systems where $\mu$ is implicit in $x$, $M(\mathcal{S})$
reduces to a chain on $X$.

We say $\mathcal{S}$ is causally emergent at scale $\Pi$ iff $\mathrm{
CE}(M(\mathcal{S}), \Pi) > 0$. We say $\mathcal{S}$ is causally
emergent iff this holds for some $\Pi$.

\subsection*{Caveats}

Three caveats matter:

\begin{itemize}[leftmargin=1.5em,itemsep=2pt]
  \item \textbf{Intervention distribution is conventional.}
    $\mathbf{H}_{\max}$ is not derived. Results may shift under other
    choices.
  \item \textbf{Sup-over-$\Pi$ is uncomputable for large $|X|$.} For
    Flow-Lenia with a $128 \times 128$ field, $|X|$ is astronomical
    and we cannot exhaustively search partitions. Tractable
    relaxations exist (greedy partition refinement,
    information-bottleneck-style optimisation) but do not exhaust the
    supremum.
  \item \textbf{Whether OP1 holds for the Flow-Lenia + LLM system is
    open.} See Chapter~9.
\end{itemize}

\section{Measurement, part two: observer-relative novelty}

We have an observer $O_s$ with state $a_t$. We have a trajectory
$(x_t, \mu_t)_{t \geq 0}$. We want a quantity that measures novelty.

For the Lenia + LLM composite, $O_{\mathrm{CLIP}}$ computes, for each
new (frame, name) pair, the cosine novelty against the archive:
\[
  \mathrm{novelty}(z_t) \;=\; \frac{1}{k} \sum_{z \in \mathrm{NN}_k(z_t,
  A)} 1 - \mathrm{cos}(z_t, z),
\]
where $\mathrm{NN}_k(z_t, A)$ is the $k$-nearest neighbours of $z_t$ in
archive $A$, and $\mathrm{cos}$ is cosine similarity. If $z_t$ is far
from everything in $A$, novelty is high; the archive update rule
adds $z_t$ to $A$ (subject to a size cap).

This is one observer family among three (Chapter~7), and one
instance of one family. The framework does not impose a single
novelty measure. It does require that the chosen one is reported
explicitly in the system description.

\subsection*{Open-endedness, operationalised}

A system exhibits observer-relative open-endedness with respect to
$O$ iff $O$'s novelty signal does not saturate over time. Three
operationalisations:

\begin{itemize}[leftmargin=1.5em,itemsep=2pt]
  \item \textbf{Hughes--Dennis learnable novelty}: the rate of
    learnable new skills is bounded away from zero in the long run.
  \item \textbf{$\Omega$ metric}: a residence-time-based measure of
    how long the system stays in any descriptor-space region. Long
    tails are open-ended.
  \item \textbf{Heaps' law novelty rate}: the cumulative count of
    distinct (frame, name) classes grows as $t^\beta$ with $0 < \beta
    < 1$.
\end{itemize}

The book takes none of these as definitive. They are three different
shapes of ``not eventually saturated.''

\section{What the slot decomposition buys you at measurement time}

A reader who has followed the book this far might ask: did we need
the slot decomposition for any of this? Could we not just measure
effective information on whatever Markov chain we happened to be
studying?

The decomposition does three things at measurement time.

\paragraph{Slot-wise ablations.} You can swap $V_{\mathrm{LLM}}$ for
random prompts and recompute $\mathrm{EI}(M(\mathcal{S}))$. The
difference attributes a portion of the system's causal emergence to
the LLM channel specifically. Without the slot decomposition, the
attribution is qualitative; with it, the experimental design follows
directly from the lens.

\paragraph{Cross-substrate comparison.} You can compare two systems
that share four slots and differ in one. Lenia + LLM vs Lenia +
random-prompts. POET vs POET-without-environment-evolution. The lens
makes the comparison principled: same slot structure, one slot
varied, score the difference.

\paragraph{Reproducibility.} The system description (see
\texttt{RESULTS.md}) requires you to record the value of every
structural slot. Two implementations are slot-wise comparable iff they
agree on the typing, record the values, preserve the iteration order,
and specify the entity-detection rule. Reproducibility is mechanical
once the lens has been adopted.

\section{Where to next}

The lens is built. The worked example is specified. The measurement
framework is in place. Chapter~9 is the frontiers chapter: the
conjectures C1--C3, the open problems OP1--OP3, the questions the
framework cannot yet answer, and the directions a reader might pursue.

The math primer at the back of the book covers the formal machinery
we have been using: information theory (entropy, mutual information,
KL divergence, effective information), measure theory (measurable
spaces, measures, Markov kernels), order theory (posets, lattices,
closure operators), and a light dose of category theory (objects,
morphisms, monoidal categories) for the $\boxtimes_\kappa$ and
$\diamond_\rho$ compositions.

\section*{To think about}

\begin{enumerate}
  \item Pick one slot of the Flow-Lenia + LLM system and propose an
    ablation. What does the ablated system tell you about the role
    of that slot?
  \item Effective information uses the uniform intervention
    distribution by convention. Suggest an alternative for a system
    whose state space has natural non-uniform structure. What changes?
  \item The $\Omega$ metric is residence-time-based. Why does long
    residence time signal open-endedness rather than the opposite?
  \item The lens treats novelty as observer-relative. Construct two
    observers for the Flow-Lenia + LLM system that would disagree
    about whether a particular run is novel.
  \item Coarse-graining search is uncomputable for large state
    spaces. What is a reasonable heuristic search procedure for
    finding emergent macroscales in a Flow-Lenia run?
\end{enumerate}
